% Three lines under the header, for the reader who will not reach page two:
% what the job is, how long, and where the evidence lives.
%
% The first clause carries the two nouns a reviewer screens on, AWS and
% Kubernetes, rather than saving them for the margin column.
%
% The site gets its own line rather than trailing the paragraph, because a
% bare domain set in body text reads as prose and nobody clicks it.
%
% \printinfo carries only the domain, so the chain glyph sits against the link
% it labels rather than leading the whole line. It is the same macro and the
% same glyph as the homepage entry in the contact line above, so the icon and
% its 0.5em gap match the header instead of being re-invented here, and the
% \mbox inside it keeps the glyph from ever being orphaned from the domain by
% a line break. The label stays outside, where it is free to break.
%
% hyperref runs with `hidelinks', so the glyph and the monospace are the only
% affordances the link gets; there is no second color on this page to spend.
% \texttt is Plex Mono, set in the preamble -- the domain reads as a token to
% be copied rather than as prose, which is what the chain glyph promises. Not
% bolded as well: mono already separates it, and both at once reads heavy.
%
% No full stop after the domain: it ends up inside the string when a reader
% copies the line.

I build the systems other engineers ship through: infrastructure as code,
GitOps delivery, and the review gates that make a change safe to approve and
possible to undo. Ten years in infrastructure, the last four owning the AWS
and Kubernetes foundation a few hundred engineers deploy through.

\smallskip
Case studies and writing:
\printinfo{\homepagesymbol}{\href{https://oponomarov.com}{\texttt{oponomarov.com}}}
